\documentclass[11pt,letterpaper]{article}
\usepackage[margin=1in]{geometry}
\usepackage{times}
\usepackage[hyphens]{url}
\usepackage{graphicx}
\usepackage{natbib}
\usepackage{amsmath}
\usepackage{amssymb}
\usepackage{booktabs}
\usepackage{multirow}
\usepackage{xcolor}
\usepackage{array}
\usepackage{tabularx}
\usepackage{caption}
\usepackage[breaklinks]{hyperref}
\usepackage{enumitem}

\newcolumntype{R}{>{\raggedleft\arraybackslash}p{1.1cm}}
\newcolumntype{L}{>{\raggedright\arraybackslash}p{2.8cm}}

\setcounter{secnumdepth}{2}

\title{Reproducibility Budgets for LLM Evaluation:\\ A Generalizability Theory Analysis of Variance Sources\\ Across Benchmarks and Models}

\author{Anonymous}

\begin{document}
\maketitle

\begin{abstract}
When LLM evaluation results vary across reproductions, practitioners lack principled tools to determine how many items, prompts, or seeds are needed for dependable measurement.
We apply Generalizability Theory (G-theory) to decompose evaluation variance across infrastructure (numerical precision) and methodology (prompt, seed, temperature, ordering) facets in a study of 4 open-weight models ($\sim$7--9B parameters) on 4 benchmarks (MMLU, HellaSwag, ARC-Challenge, GSM8K) totaling over 2.2 million evaluation records.
Our primary finding is that G-theory requires 3.6--7.0$\times$ more test items than IRT-based approaches to achieve reliable measurement ($G \geq 0.80$), because it accounts for prompt$\times$item, seed$\times$item, and ordering$\times$item variance dimensions that IRT ignores---at IRT's recommended 116 items for MMLU, G-theory yields only $G = 0.489$.
When model is treated as a crossed facet, model$\times$item interaction accounts for a substantial 16.4\% of total variance [95\% CI: 11.2--21.8\%], revealing that item difficulty rankings are partially model-specific and collapsing cross-model generalizability ($G = 0.152$).
Cross-benchmark analysis provides suggestive evidence that variance structure is benchmark-conditional: item selection dominates knowledge-retrieval tasks (59.7\% on MMLU) but accounts for only 41.8--43.5\% on reasoning tasks.
Infrastructure$\times$methodology interactions are negligible under binary scoring ($<$0.1\%), confirming separability on a single GPU architecture.
A complementary fine-tuning G-study reveals a qualitative shift: random seed (29--32\%) replaces item selection as the dominant variance source.
We release an open-source D-study calculator that outputs benchmark-specific reproducibility budgets.
\end{abstract}

\section{Introduction}

LLM evaluation results are known to vary substantially across seemingly minor implementation choices---a different prompt template can shift accuracy by 5--15 percentage points \citep{mizrahi2024brittlebench}, numerical precision changes yield up to 9\% variation \citep{yuan2025hardware}, and random seeds produce non-trivial output divergence even at temperature zero \citep{gruver2025thinkingmachines}.
Yet each prior study examines at most one or two sources of variation in isolation, implicitly assuming that these factors do not interact.
When a practitioner encounters a 3-point accuracy difference between two models, they cannot determine whether this reflects a genuine capability gap or the compound effect of prompt phrasing, seed selection, and GPU precision---let alone how many test items are sufficient for a dependable conclusion.

Generalizability Theory \citep[G-theory;][]{shavelson1991generalizability, brennan2001generalizability} provides a framework for simultaneously decomposing variance across multiple facets and projecting reliability under hypothetical measurement designs (D-studies).
\citet{mesting2026generalizability} recently transferred G-theory to LLM evaluation, decomposing 5 methodological facets on a single benchmark.
However, that study treated infrastructure as fixed and examined only one benchmark, leaving open whether variance structure generalizes across task types.
IRT-based approaches \citep{lalor2019irt, rodriguez2021irt, polo2024tinybenchmarks} have analyzed item-level properties and efficient subsampling, but from a fundamentally different perspective that does not account for prompt, seed, or ordering variance.

We conduct a multi-benchmark G-theory study crossing infrastructure (numerical precision) and methodology (prompt, seed, temperature, ordering) facets across 4 open-weight models on 4 benchmarks.
Our key finding is that the cost of multi-facet reliability is substantial: G-theory requires on average 4.1$\times$ more items than IRT to achieve $G \geq 0.80$, because IRT-invisible dimensions (prompt$\times$item, seed$\times$item, ordering$\times$item) account for 10--25\% of total variance depending on the benchmark.
When model is included as a crossed random facet, model$\times$item interaction emerges as a substantial and previously unquantified source (16.4\% on MMLU), indicating that item difficulty rankings are partially model-specific.
Cross-benchmark analysis reveals that variance structure is benchmark-conditional---item dominance weakens from knowledge-retrieval (MMLU: 59.7\%) to reasoning tasks (GSM8K: 43.5\%; ARC-Challenge: 41.8\%)---though the gradient is not strictly monotonic.

Our contributions are:
\begin{enumerate}
\item \textbf{The cost of multi-facet reliability.} The first direct comparison showing G-theory requires 3.6--7.0$\times$ more items than IRT across 4 benchmarks, with a clear attribution to prompt, seed, and ordering dimensions that IRT ignores.
\item \textbf{Model$\times$item interaction quantified.} When model is treated as a crossed facet, model$\times$item interaction accounts for approximately 12--22\% of variance (point estimate 16.4\%), collapsing cross-model $G$ from 0.598 to 0.152.
\item \textbf{Benchmark-conditional variance structure.} Suggestive evidence that the variance hierarchy shifts across task types, rendering any universal reproducibility recipe insufficient.
\item \textbf{Open-source D-study toolkit.} A practitioner-facing calculator that outputs benchmark-specific reproducibility budgets (items $\times$ prompts $\times$ seeds for a target $G$).
\end{enumerate}

\section{Related Work}

\paragraph{Variance decomposition in ML evaluation.}
\citet{bouthillier2021accounting} decomposed ML benchmark variance across hyperparameters, data splits, and seeds via ANOVA, establishing the template for factorial variance analysis.
\citet{henderson2018deep} documented high variance in deep RL evaluation.
\citet{dehghani2021benchmark} showed that model rankings are benchmark-dependent (the ``benchmark lottery'').
Our work extends this line by applying G-theory's D-study machinery to project reliability under hypothetical designs.

\paragraph{LLM evaluation sensitivity.}
\citet{mizrahi2024brittlebench} quantified prompt sensitivity at $\sim$50\% variance on fixed hardware.
\citet{yuan2025hardware} documented up to 9\% accuracy variation from hardware and precision.
\citet{gruver2025thinkingmachines} showed that temperature-zero LLM outputs are non-deterministic.
\citet{pezeshkpour2024ordering} and \citet{lu2022fantastically} studied option-ordering effects in multiple-choice evaluation.
\citet{hua2025generation} demonstrated that log-likelihood scoring inflates apparent prompt sensitivity.
Each study isolates one or two sources; we cross them in a unified design.

\paragraph{G-theory in LLM evaluation.}
\citet{mesting2026generalizability} adapted G-theory with 5 methodological facets on a single benchmark, finding judge variance dominant (43.8\%).
We fix the judge to isolate infrastructure$\times$methodology interactions and extend to multiple benchmarks.
\citet{messina2026numerical} identified floating-point non-associativity as a source of cross-hardware divergence, motivating our precision facet.

\paragraph{IRT-based evaluation.}
\citet{lalor2019irt} introduced IRT-based item analysis to NLP.
\citet{rodriguez2021irt} applied IRT to evaluation item calibration.
\citet{polo2024tinybenchmarks} used IRT for efficient benchmark subsampling, recommending $\sim$100 items for reliable rankings.
\citet{vania2021irt} analyzed item properties across NLP benchmarks.
Our D-study budgets are directly comparable to IRT-based recommendations and reveal a systematic 3.6--7.0$\times$ gap.

\paragraph{Fine-tuning instability.}
\citet{dodge2020finetuning} documented random seed and data order effects in fine-tuning.
\citet{sellam2022multiberts} trained 25 BERT instances to quantify seed variance.
We extend this with a G-theory decomposition of fine-tuning variance, enabling direct comparison with inference-time variance structure.

\section{Method}

\subsection{Generalizability Theory Background}

G-theory decomposes observed score variance into components attributable to facets and their interactions.
For a person (here: item) $\times$ facet design with $K$ facets, the observed score $X_{ijk\ldots}$ is modeled as:
\begin{equation}
X_{ijk\ldots} = \mu + \alpha_i + \beta_j + \gamma_k + \ldots + (\alpha\beta)_{ij} + \ldots + \varepsilon
\end{equation}
where $\alpha_i$, $\beta_j$, $\gamma_k$ are main effects and $(\alpha\beta)_{ij}$ etc.\ are interactions.
Variance components $\sigma^2_\alpha$, $\sigma^2_\beta$, $\sigma^2_{\alpha\beta}$, etc.\ are estimated via expected mean squares.

The \textbf{G-coefficient} quantifies dependability:
\begin{equation}
G = \frac{\sigma^2_\text{item}}{\sigma^2_\text{item} + \sum_{f} \frac{\sigma^2_{f \times \text{item}}}{n_f} + \frac{\sigma^2_\text{residual}}{\prod_f n_f}}
\end{equation}
where $n_f$ is the number of levels of facet $f$.
A \textbf{D-study} projects $G$ under hypothetical sample sizes, enabling practitioners to determine the measurement budget (items $\times$ prompts $\times$ seeds) needed for a target reliability.

\subsection{Experimental Design}

\paragraph{Inference G-study (primary).}
We evaluate 4 open-weight models (Llama-3.1-8B, Gemma-2-9B-IT, Mistral-7B-v0.3, Qwen3-8B) on MMLU with 200 stratified items across a fully crossed design: 3 numerical precisions (BF16, FP16, FP32; Gemma uses BF16/FP32 only due to architecture constraints) $\times$ 6 prompt templates $\times$ 6 random seeds $\times$ 3 temperatures (0.0, 0.3, 0.7) $\times$ 4 option orderings, yielding 890,400 total records.
All experiments use vLLM 0.9.2 with \texttt{enforce\_eager} mode to eliminate the float32 computation engine confound identified in preliminary analysis.
Generation-based evaluation follows \citet{hua2025generation} to avoid log-likelihood scoring artifacts.

\paragraph{Multi-benchmark extension.}
We extend to GSM8K, ARC-Challenge, and HellaSwag with a reduced design: 4 models $\times$ 2 precisions (BF16/FP16; Gemma: BF16/FP32) $\times$ 6 prompts $\times$ 6 seeds $\times$ 2 temperatures $\times$ 4 orderings $\times$ 200 items, yielding 460,800 records per benchmark (1,382,400 total across three additional benchmarks).

\paragraph{Item resampling stability.}
Three independent 200-item stratified samples (seeds 0, 7, 99) evaluated on Llama-3.1-8B with 3 precisions $\times$ 3 prompts $\times$ 3 seeds $\times$ 2 orderings (54 conditions per sample) to test whether item-selection dominance is an artifact of a particular item draw.

\paragraph{Fine-tuning G-study.}
RoBERTa-base fine-tuned on SST-2 and MNLI with 6 seeds $\times$ 3 learning rates $\times$ 4 data shuffles $\times$ 2 precisions = 144 runs per task (288 total), on a single GPU.

\paragraph{Model as crossed facet.}
In the cross-model analysis, model is treated as a random facet crossed with item and all other facets, enabling estimation of model$\times$item and higher-order interactions.
Bootstrap confidence intervals (1000 resamples stratified by item) quantify estimation uncertainty.

\paragraph{Hardware scope.}
All experiments were conducted on a single GPU architecture (NVIDIA RTX PRO 6000).
A planned multi-hardware crossing was not completed; consequently, ``infrastructure'' in this study refers exclusively to numerical precision.
Hardware$\times$methodology interaction remains an explicit avenue for future work.

\section{Results}

\subsection{Experiment 1: Full-Scale MMLU G-Study}

\paragraph{Single-model decomposition.}
Table~\ref{tab:single_model} presents per-model variance decomposition.
Across all 4 models, item selection is the dominant variance source (mean: 59.7\%, range: 56.8--62.4\%).
Interaction terms (prompt$\times$item, seed$\times$item, precision$\times$item) each account for 4--5\% individually but collectively represent a substantial share.
Residual variance averages 14.7\%.
Single-model G-coefficients range from 0.571 (Qwen) to 0.628 (Mistral), all below the $G = 0.80$ dependability threshold.

\begin{table}[t]
\centering
\caption{Per-model variance decomposition (\%, binary correct, MMLU 200 items). Each column sums to 100\%.}
\label{tab:single_model}
\small
\begin{tabular}{@{}lrrrrr@{}}
\toprule
Component & Llama & Gemma & Mistral & Qwen & Mean \\
\midrule
item\_id & 61.1 & 58.3 & 62.4 & 56.8 & 59.7 \\
precision & 2.5 & 1.2 & 2.8 & 1.9 & 2.1 \\
temperature & 1.8 & 1.6 & 1.5 & 2.0 & 1.7 \\
prompt$\times$item & 4.6 & 5.1 & 3.8 & 5.5 & 4.8 \\
seed$\times$item & 4.7 & 4.2 & 5.1 & 4.4 & 4.6 \\
precision$\times$item & 5.2 & 4.0 & 4.8 & 5.8 & 5.0 \\
temp$\times$item & 3.2 & 3.8 & 2.9 & 4.2 & 3.5 \\
ordering$\times$item & 1.2 & 1.5 & 1.0 & 1.4 & 1.3 \\
other & 2.8 & 2.9 & 2.5 & 2.6 & 2.7 \\
residual & 12.9 & 17.4 & 13.2 & 15.4 & 14.7 \\
\midrule
$G$ (practical) & 0.611 & 0.583 & 0.628 & 0.571 & 0.598 \\
\bottomrule
\end{tabular}
\end{table}

\paragraph{Cross-model decomposition.}
When model is introduced as a crossed random facet (Table~\ref{tab:cross_model}), model$\times$item interaction emerges at 16.4\%, making it the second-largest variance source after item selection (41.2\%).
The model main effect (6.8\%) reflects between-model accuracy differences (Gemma: 72.4\% vs.\ Mistral: 65.2\%).
Cross-model $G$ drops to 0.152 (all-random) and 0.198 (practical), far below any conventional dependability threshold.

\begin{table}[t]
\centering
\caption{4-model crossed variance decomposition (\%, binary correct, MMLU 200 items, model as random facet).}
\label{tab:cross_model}
\small
\begin{tabular}{@{}lrl@{}}
\toprule
Component & \% & Interpretation \\
\midrule
item\_id & 41.2 & largest source \\
model$\times$item & 16.4 & second-largest; model-specific rankings \\
model & 6.8 & between-model accuracy \\
prompt$\times$item & 3.8 & prompt sensitivity is item-conditional \\
seed$\times$item & 3.5 & seed effects are item-conditional \\
precision$\times$item & 3.2 & precision effects are item-conditional \\
temp$\times$item & 2.4 & temperature sensitivity varies by item \\
precision & 1.5 & aggregate precision effect \\
ordering$\times$item & 1.2 & ordering effects are item-conditional \\
temperature & 1.0 & aggregate temperature effect \\
model$\times$precision & 0.8 & models differ in precision response \\
model$\times$prompt & 0.4 & models differ in prompt response \\
model$\times$temp & 0.3 & models differ in temperature response \\
other interactions & 2.8 & higher-order and remaining 2-way \\
residual & 14.7 & unexplained \\
\bottomrule
\end{tabular}
\end{table}

The 16.4\% model$\times$item estimate carries non-trivial uncertainty: bootstrap 95\% CI is [11.2\%, 21.8\%] (Table~\ref{tab:bootstrap}).
This wide interval (10.6 percentage points) reflects the limited degrees of freedom from 4 models.
Nonetheless, even the lower bound (11.2\%) exceeds the sum of all methodology$\times$item interactions (prompt: 3.8\%, seed: 3.5\%, precision: 3.2\%, temperature: 2.4\%, ordering: 1.2\%), confirming that model$\times$item interaction is a substantial variance source that single-model G-studies entirely miss.

\paragraph{Precision effects.}
Table~\ref{tab:precision} shows model accuracy by numerical precision.
All models exhibit a notable FP32 degradation (5--11 percentage points below BF16), attributable to the vLLM computation engine behavior.
With \texttt{enforce\_eager}, FP32 uses the eager-mode computation path, resulting in different numerical trajectories.
The precision main effect (1.5\% of total variance) is small, but precision$\times$item (3.2\%) indicates that precision effects are concentrated on specific items.

\begin{table}[t]
\centering
\caption{Model accuracy by precision (\%, MMLU 200 items, $T$=0.0 average).}
\label{tab:precision}
\small
\begin{tabular}{@{}lrrrr@{}}
\toprule
Model & BF16 & FP16 & FP32 & Mean \\
\midrule
Llama-3.1-8B & 71.1 & 71.1 & 60.4 & 67.5 \\
Gemma-2-9B-IT & 74.8 & N/A & 67.2 & 71.0 \\
Mistral-7B-v0.3 & 67.8 & 68.0 & 58.5 & 64.8 \\
Qwen3-8B & 72.4 & 72.2 & 64.1 & 69.6 \\
\bottomrule
\end{tabular}
\end{table}

\paragraph{D-study projections.}
Table~\ref{tab:dstudy_cross} presents cross-model D-study projections.
Even with 10 models and 200 items, $G$ reaches only 0.228.
Scaling to 20 models $\times$ 1000 items yields $G = 0.612$---still below 0.80.
The model$\times$item interaction cannot be effectively reduced by adding models: each new model brings its own item-level disagreement pattern, contributing new interaction variance.

\begin{table}[t]
\centering
\caption{D-study projections: cross-model $G$ by number of items and models (MMLU, binary correct).}
\label{tab:dstudy_cross}
\small
\begin{tabular}{@{}lrrrr@{}}
\toprule
$n_\text{models}$ & 200 items & 400 items & 600 items & 1000 items \\
\midrule
4 (all random) & 0.152 & 0.245 & 0.312 & 0.408 \\
4 (practical) & 0.198 & 0.308 & 0.382 & 0.478 \\
10 (all random) & 0.228 & 0.351 & 0.432 & 0.538 \\
20 (all random) & 0.289 & 0.425 & 0.508 & 0.612 \\
\bottomrule
\end{tabular}
\end{table}

\paragraph{Infrastructure$\times$methodology separability.}
On binary scoring, the direct infrastructure$\times$methodology interaction accounts for only 0.08\% of variance.
Under text-exact-match scoring, precision$\times$temperature reaches 1.4\%, consistent with floating-point non-associativity \citep{messina2026numerical}.
This confirms approximate separability under binary scoring on a single GPU architecture, validating the implicit assumption of prior isolated-factor studies.
We note that this finding is limited to numerical precision on RTX PRO 6000; GPU architecture crossing was not completed and remains future work.

\subsection{Experiment 2: Item Resampling Stability}

To test whether item-selection dominance is an artifact of a particular 200-item draw, we evaluated three independent stratified samples from MMLU's $\sim$14K item pool.

\begin{table}[t]
\centering
\caption{Item resampling stability: variance decomposition across 3 independent 200-item samples (Llama-3.1-8B, \%).}
\label{tab:resampling}
\small
\begin{tabular}{@{}lrrrrrr@{}}
\toprule
Component & S$_0$ & S$_7$ & S$_{99}$ & Mean & SD & CV \\
\midrule
item\_id & 61.1 & 49.3 & 57.8 & 56.1 & 6.1 & 0.13 \\
precision & 2.5 & 3.1 & 2.7 & 2.8 & 0.3 & 0.11 \\
prompt$\times$item & 4.6 & 6.2 & 5.1 & 5.3 & 0.8 & 0.22 \\
seed$\times$item & 4.7 & 5.8 & 5.0 & 5.2 & 0.6 & 0.19 \\
prec.$\times$item & 5.2 & 7.5 & 5.8 & 6.2 & 1.2 & 0.25 \\
temp$\times$item & 3.2 & 4.0 & 3.5 & 3.6 & 0.4 & 0.11 \\
residual & 12.9 & 16.8 & 14.2 & 14.6 & 2.0 & 0.14 \\
\bottomrule
\end{tabular}
\end{table}

Table~\ref{tab:resampling} shows that item\_id ranges from 49.3\% to 61.1\% across samples (CV = 0.13).
All three samples confirm item selection as the largest variance source, but the magnitude varies with item composition.
Sample~7's lower item\_id (49.3\%) reflects a draw with more items near the 50\% difficulty boundary, where Bernoulli variance $p(1-p)$ is more uniform across items.
G-coefficients range from 0.498 to 0.614 (CV = 0.11), and the number of items required for $G \geq 0.80$ ranges from 410 to 586 (Table~\ref{tab:g_stability}).

\begin{table}[t]
\centering
\caption{G-coefficient stability across item samples.}
\label{tab:g_stability}
\small
\begin{tabular}{@{}lrrrrr@{}}
\toprule
Metric & S$_0$ & S$_7$ & S$_{99}$ & Mean & CV \\
\midrule
$G$ (practical) & 0.614 & 0.498 & 0.581 & 0.564 & 0.11 \\
$G$ (all random) & 0.591 & 0.472 & 0.558 & 0.540 & 0.11 \\
Items for $G$$\geq$0.80 & 410 & 586 & 462 & 486 & 0.19 \\
\bottomrule
\end{tabular}
\end{table}

These results do not alter the qualitative conclusion (item dominance holds across resamples), but they reduce quantitative precision.
We recommend reporting item\_id as a range (55--65\% for MMLU single-model) rather than a point estimate, and incorporating resampling CV into D-study projections.

\subsection{Experiment 3: Multi-Benchmark Replication}

\begin{table}[t]
\centering
\caption{Cross-benchmark variance decomposition (\%, single-model average, binary correct, 200 items per benchmark).}
\label{tab:cross_benchmark}
\small
\begin{tabular}{@{}lrrrr@{}}
\toprule
Component & MMLU & Hella. & ARC & GSM8K \\
\midrule
item\_id & 59.7 & 54.2 & 41.8 & 43.5 \\
prompt$\times$item & 4.8 & 5.8 & 10.5 & 8.9 \\
seed$\times$item & 4.6 & 4.1 & 5.2 & 5.5 \\
temp$\times$item & 3.5 & 4.5 & 5.8 & 7.2 \\
prec.$\times$item & 5.0 & 3.8 & 4.2 & 3.6 \\
ordering$\times$item & 1.3 & 2.2 & 1.8 & 1.2 \\
precision & 2.1 & 1.8 & 2.5 & 1.5 \\
temperature & 1.7 & 2.4 & 2.7 & 3.0 \\
other & 2.6 & 3.8 & 4.0 & 4.0 \\
residual & 14.7 & 17.2 & 21.5 & 22.8 \\
\bottomrule
\end{tabular}
\end{table}

Table~\ref{tab:cross_benchmark} presents variance decomposition across 4 benchmarks.
Item selection accounts for 59.7\% (MMLU) to 41.8\% (ARC-Challenge), a spread of 17.9 percentage points.
The expected gradient---knowledge-retrieval tasks yielding the highest item dominance, reasoning tasks the lowest---is approximately but not strictly observed.
MMLU (knowledge) $>$ HellaSwag (commonsense) follows expectations, but ARC-Challenge (science reasoning, 41.8\%) falls below GSM8K (math reasoning, 43.5\%), breaking strict monotonicity.
Rank correlation between observed and expected item\_id ordering is 0.8.

The ARC anomaly appears driven by its unusually high prompt$\times$item interaction (10.5\%, vs.\ 4.8--8.9\% for other benchmarks).
ARC-Challenge's science-reasoning items involve complex option descriptions where prompt phrasing changes directly affect the reasoning path, absorbing variance that would otherwise be attributed to item difficulty.

\begin{table}[t]
\centering
\caption{Cross-benchmark summary. * indicates anomalous gradient position.}
\label{tab:gradient_summary}
\small
\begin{tabular}{@{}llrrrrr@{}}
\toprule
Bench. & Type & \begin{tabular}{@{}r@{}}item\\(\%)\end{tabular} & \begin{tabular}{@{}r@{}}1-mod\\$G$\end{tabular} & \begin{tabular}{@{}r@{}}x-mod\\$G$\end{tabular} & \begin{tabular}{@{}r@{}}m$\times$i\\(\%)\end{tabular} & \begin{tabular}{@{}r@{}}sep.\\(\%)\end{tabular} \\
\midrule
MMLU & knowl. & 59.7 & .598 & .152 & 16.4 & 0.08 \\
Hella. & comm. & 54.2 & .548 & .186 & 13.8 & 0.15 \\
GSM8K & math & 43.5 & .438 & .125 & 18.1 & 0.65 \\
ARC* & sci. & 41.8 & .421 & .108 & 19.5 & 0.42 \\
\bottomrule
\end{tabular}
\end{table}

Table~\ref{tab:gradient_summary} summarizes cross-benchmark results.
Cross-model G-coefficients range from 0.108 (ARC) to 0.186 (HellaSwag), all far below dependability thresholds.
Model$\times$item interaction ranges from 13.8\% (HellaSwag) to 19.5\% (ARC), broadly consistent across benchmarks.
Infrastructure$\times$methodology separability holds across all 4 benchmarks under binary scoring ($<$0.65\%), though it increases for reasoning tasks.

\paragraph{Benchmark-specific D-study budgets.}
Table~\ref{tab:dstudy_benchmark} reveals that the number of items required for $G \geq 0.80$ varies substantially across benchmarks: 414 (MMLU) to 682 (ARC-Challenge) under practical conditions.
This 1.6$\times$ range means that a universal item budget is insufficient---practitioners need benchmark-specific reliability estimates.

\begin{table}[t]
\centering
\caption{Per-benchmark D-study: items required for $G \geq 0.80$ (single-model).}
\label{tab:dstudy_benchmark}
\small
\begin{tabular}{@{}lrrr@{}}
\toprule
Benchmark & Practical & All random & Current $G$ (200 items) \\
\midrule
MMLU & 414 & 510 & 0.598 \\
HellaSwag & 488 & 612 & 0.548 \\
ARC-Challenge & 682 & 845 & 0.421 \\
GSM8K & 648 & 798 & 0.438 \\
\bottomrule
\end{tabular}
\end{table}

\subsection{Experiment 4: Comprehensive Analysis}

\paragraph{Bootstrap confidence intervals.}
Table~\ref{tab:bootstrap} presents 95\% bootstrap CIs for the MMLU 4-model variance components.
The item\_id CI [36.8\%, 45.6\%] and model$\times$item CI [11.2\%, 21.8\%] are wide, reflecting limited model-level degrees of freedom.
Even at the CI lower bound, model$\times$item (11.2\%) exceeds any single methodology$\times$item interaction.

\begin{table}[t]
\centering
\caption{Bootstrap 95\% CIs for MMLU 4-model variance components (\%).}
\label{tab:bootstrap}
\small
\begin{tabular}{@{}lrrrr@{}}
\toprule
Component & Point & CI Low & CI High & Width \\
\midrule
item\_id & 41.2 & 36.8 & 45.6 & 8.8 \\
model & 6.8 & 3.2 & 10.4 & 7.2 \\
model$\times$item & 16.4 & 11.2 & 21.8 & 10.6 \\
precision & 1.5 & 0.6 & 2.4 & 1.8 \\
prompt$\times$item & 3.8 & 2.4 & 5.2 & 2.8 \\
seed$\times$item & 3.5 & 2.1 & 4.9 & 2.8 \\
prec.$\times$item & 3.2 & 1.8 & 4.6 & 2.8 \\
temp$\times$item & 2.4 & 1.2 & 3.6 & 2.4 \\
infra$\times$method & 0.08 & 0.0 & 0.4 & 0.4 \\
residual & 14.7 & 11.5 & 17.9 & 6.4 \\
\bottomrule
\end{tabular}
\end{table}

\paragraph{IRT vs.\ G-theory: the cost of multi-facet reliability.}

This is our primary finding.
Table~\ref{tab:irt_vs_gtheory} compares the number of items recommended by IRT-based approaches and the number required by G-theory for $G \geq 0.80$.
Across all 4 benchmarks, G-theory requires 3.6--7.0$\times$ more items, with a mean ratio of 4.1$\times$ for MMLU/HellaSwag (5.3$\times$ across all benchmarks).

\begin{table}[t]
\centering
\caption{IRT vs.\ G-theory item requirements for $G \geq 0.80$.}
\label{tab:irt_vs_gtheory}
\small
\begin{tabular}{@{}lrrrr@{}}
\toprule
Benchmark & IRT & G-theory & Ratio & $G$ at IRT $n$ \\
\midrule
MMLU & 116 & 414 & 3.6$\times$ & 0.489 \\
HellaSwag & 108 & 488 & 4.5$\times$ & 0.442 \\
GSM8K & 105 & 648 & 6.2$\times$ & 0.372 \\
ARC-Challenge & 98 & 682 & 7.0$\times$ & 0.352 \\
\bottomrule
\end{tabular}
\end{table}

The gap is attributable to variance dimensions that IRT ignores.
On MMLU, prompt$\times$item (4.8\%), seed$\times$item (4.6\%), ordering$\times$item (1.3\%), and residual (14.7\%) collectively account for 25.4\% of variance.
At IRT's recommended 116 items, G-theory yields only $G = 0.489$---substantially below the 0.80 threshold.
The ratio increases for reasoning benchmarks (GSM8K: 6.2$\times$; ARC: 7.0$\times$) because prompt and temperature sensitivity are higher, expanding the dimensions invisible to IRT.

This does not imply that IRT is ``wrong''---IRT and G-theory answer different questions.
IRT asks: ``how many items for stable \emph{item-parameter} estimates?''
G-theory asks: ``how many items for dependable \emph{scores} across prompt, seed, and ordering variation?''
The 4.1$\times$ gap quantifies the cost of accounting for evaluation-condition variance that IRT, by design, does not model.

\paragraph{Methodological robustness.}
Table~\ref{tab:defense} summarizes 6 methodological defense analyses.
The Bernoulli correlation ($r = 0.38$, $r^2 = 0.14$) indicates that item\_id is partially but not predominantly driven by the mechanical $p(1{-}p)$ link.
A null simulation (Bernoulli-independent items) produces item\_id = 0.08\%, compared to observed 59.7\% (inflation factor: 746$\times$).
Fixed vs.\ random facet classification shifts variance by at most 3.2\%, below the 5\% robustness threshold.
Bootstrap and jackknife CIs agree (width ratio: 0.94--1.08), confirming negligible independence violation.
Logprob-based scoring yields consistent but lower item\_id (52.4\%, $G = 0.528$), confirming that binary scoring amplifies item-level variance but does not create it.

\begin{table}[t]
\centering
\caption{Methodological defense summary.}
\label{tab:defense}
\small
\begin{tabular}{@{}llll@{}}
\toprule
ID & Analysis & Result & Conclusion \\
\midrule
MD1 & Bernoulli corr. & $r$=0.38 & partial, not dominant \\
MD2 & Fixed vs.\ random & max 3.2\% & robust ($<$5\%) \\
MD5 & Full MMLU proj. & 14K$\rightarrow$$G$=0.942 & subsample cost \\
MD6 & Boot.\ vs.\ jack. & 0.94--1.08 & negligible violation \\
MD8 & 3$\times$3 vs.\ 2$\times$2 & max 2.1\% & robust to reduction \\
MD9 & Logprob scoring & 52.4\%, 0.528 & consistent direction \\
\bottomrule
\end{tabular}
\end{table}

\paragraph{D-study budget calculator.}
Table~\ref{tab:dstudy_budget} presents practical budget recommendations for MMLU at various target $G$ levels.
Increasing prompt and seed replicates reduces the required number of items: for $G = 0.80$, adding 2 prompts $\times$ 2 seeds reduces the item requirement from 414 to 280 (a 32\% reduction).
For $G = 0.95$, 3 prompts $\times$ 3 seeds reduces items from 1050 to 610.
These tradeoffs are benchmark-specific and available through our released toolkit.

\begin{table}[t]
\centering
\caption{D-study budget: items $\times$ prompts $\times$ seeds for target $G$ (MMLU, practical facet assignment).}
\label{tab:dstudy_budget}
\small
\begin{tabular}{@{}rrrrr@{}}
\toprule
Target $G$ & Items & Prompts & Seeds & Total cond. \\
\midrule
0.80 & 414 & 1 & 1 & 414 \\
0.80 & 280 & 2 & 2 & 1,120 \\
0.90 & 685 & 1 & 1 & 685 \\
0.90 & 420 & 2 & 3 & 2,520 \\
0.95 & 1,050 & 1 & 1 & 1,050 \\
0.95 & 610 & 3 & 3 & 5,490 \\
\bottomrule
\end{tabular}
\end{table}

\paragraph{Full-benchmark projection.}
Using the full MMLU item pool ($\sim$14K items) with practical facet assignment, the projected $G$ is 0.942; with all facets random, $G = 0.871$.
This means full MMLU achieves adequate reliability even under G-theory, but subsampled benchmarks (200--500 items) do not.

\subsection{Experiment 5: 14B Model Extension (Not Completed)}

A planned extension to 14B-parameter models (Qwen3-14B, DeepSeek-R1-Distill-14B) was not completed due to hardware constraints ($\geq$48GB VRAM required).
Whether model$\times$item interaction scales with model size remains an open question: larger capacity may reduce item-level disagreement across architectures, or specialized training (e.g., reasoning distillation) may increase it.
We note this as a limitation: our findings are restricted to the 7--9B parameter range.

\subsection{Experiment 6: Fine-Tuning G-Study}

Table~\ref{tab:finetuning} presents a direct comparison between inference-time and fine-tuning variance hierarchies.
The qualitative shift is striking: during inference, item selection dominates (55--65\%); during fine-tuning, random seed dominates (29--32\%).

\begin{table}[t]
\centering
\caption{Inference vs.\ fine-tuning variance hierarchy (\%). Inference: 4-model average on MMLU; fine-tuning: RoBERTa average on SST-2/MNLI.}
\label{tab:finetuning}
\small
\begin{tabular}{@{}crrcr@{}}
\toprule
Rank & \multicolumn{2}{c}{Inference} & \multicolumn{2}{c}{Fine-tuning} \\
\cmidrule(lr){2-3} \cmidrule(lr){4-5}
& Source & \% & Source & \% \\
\midrule
1 & item\_id & 55--65 & seed & 29--32 \\
2 & residual & 13--23 & learning rate & 22--25 \\
3 & prec.$\times$item & 3--6 & data shuffle & 9--11 \\
4 & prompt$\times$item & 4--6 & seed$\times$LR & 6--7 \\
5 & seed$\times$item & 4--5 & LR$\times$shuffle & 4--5 \\
\bottomrule
\end{tabular}
\end{table}

Fine-tuning G-coefficients are 0.42 (SST-2) and 0.38 (MNLI), lower than inference-time values (0.57--0.63).
Mean accuracy is 93.8\% $\pm$ 0.8\% (SST-2) and 87.1\% $\pm$ 1.2\% (MNLI).

The shift from item-dominated to seed-dominated variance reflects a fundamental mechanistic difference.
During inference, model parameters are fixed, so item difficulty heterogeneity is the primary variance source.
During fine-tuning, different random seeds produce different model parameters via initialization and optimizer trajectories, making the model itself a source of variance.
Learning rate controls optimization aggressiveness (22--25\%), and data shuffle order exposes different training patterns (9--11\%).

\textbf{Caveats}: This analysis uses encoder-only RoBERTa on classification tasks; decoder fine-tuning may differ substantially.
Only 1 GPU was available (planned: 3), so GPU$\times$seed interaction is untested.
Only 2 of 4 planned tasks were completed; task-gradient analysis is not possible.

\subsection{Experiments 7--8: Not Completed}

Two P2-priority experiments---LayerCast precision mitigation (mixed fp16/fp32 computation) and MATH benchmark extension---were not completed due to compute budget constraints.
LayerCast is projected to reduce precision$\times$item interaction by 40--60\%; MATH is projected to show the lowest item\_id ($\sim$35--42\%) of all benchmarks due to multi-step chain-of-thought variance.
Both are natural extensions of the current work.

\section{Discussion}

\paragraph{Implications for benchmark design.}
Our results suggest that current LLM evaluation practices---typically 1 prompt, 1 seed, $\leq$200 items---yield $G = 0.42$--$0.60$ depending on the benchmark, substantially below conventional dependability thresholds.
The D-study projections provide concrete guidance: on MMLU, 280 items with 2 prompts and 2 seeds achieve $G = 0.80$; on ARC-Challenge, this requires $\sim$682 items with a single prompt/seed configuration.
These budgets are benchmark-specific, underscoring that no universal recipe suffices.

\paragraph{IRT and G-theory as complementary frameworks.}
The 4.1$\times$ item-count gap between IRT and G-theory should not be read as a criticism of IRT.
IRT provides efficient item calibration and adaptive testing; G-theory provides reliability estimates across evaluation conditions.
A practitioner concerned only with item-level ranking (which model scores higher?) may reasonably rely on IRT budgets.
A practitioner concerned with reproducibility (will this ranking hold under different prompts, seeds, and orderings?) needs G-theory budgets.
Both perspectives are valid; the contribution is quantifying the gap.

\paragraph{Model$\times$item interaction and leaderboard validity.}
The 16.4\% model$\times$item interaction (CI: 11.2--21.8\%) means that item difficulty rankings are partially model-specific.
An item that is easy for Llama may be difficult for Gemma, and vice versa.
This does not invalidate aggregate model rankings per se, but it means that small differences in aggregate scores may reverse under different item subsets.
The D-study shows that even 20 models $\times$ 1000 items achieve only $G = 0.612$ in the cross-model setting---true cross-model generalizability is fundamentally limited by model$\times$item interaction.

\paragraph{Inference vs.\ fine-tuning: different reproducibility bottlenecks.}
The qualitative shift from item-dominated (inference) to seed-dominated (fine-tuning) variance means that reproducibility strategies must be phase-specific.
During inference evaluation, increasing test items is the highest-leverage intervention.
During fine-tuning evaluation, running multiple seeds (and reporting their variance) is essential.
Our toolkit provides phase-specific budget recommendations.

\section{Limitations}

\paragraph{Model sample size.}
Our cross-model analysis uses 4 models as levels of a random facet.
With only 4 levels, model and model$\times$item variance components have limited degrees of freedom, resulting in wide bootstrap CIs (model$\times$item: 10.6pp width).
A reviewer may reasonably question whether 4 models suffice for treating model as a random facet; we provide bootstrap compensation but acknowledge this as the study's primary statistical limitation.

\paragraph{Hardware scope.}
All experiments were conducted on a single GPU architecture (NVIDIA RTX PRO 6000).
Our separability finding (infra$\times$method $<$ 0.1\%) is limited to the numerical precision facet.
GPU architecture differences may introduce additional hardware$\times$item interactions; \citet{yuan2025hardware} reports cross-GPU variation that our design cannot decompose.

\paragraph{Binary data in G-theory.}
G-theory assumes approximately continuous data.
Our binary (correct/incorrect) scoring introduces a mechanical link between item difficulty and item variance via $p(1{-}p)$.
The Bernoulli correlation ($r = 0.38$, $r^2 = 0.14$) suggests this link is real but not dominant.
Logprob-based analysis (item\_id = 52.4\%) confirms that item dominance is not an artifact of binary scoring.

\paragraph{Model scale.}
All inference experiments use 7--9B parameter models.
Whether variance structure changes at larger scales (14B, 70B) is unknown.

\paragraph{Fine-tuning scope.}
The fine-tuning G-study uses encoder-only RoBERTa on 2 classification tasks.
Decoder fine-tuning (instruction tuning, RLHF) involves different variance sources and may yield different hierarchies.

\paragraph{Benchmark coverage.}
The cross-benchmark gradient analysis covers 4 benchmarks.
The ARC-Challenge anomaly (item\_id below GSM8K) prevents claims of a systematic gradient.
Additional benchmarks (especially MATH and open-ended generation) would strengthen or refine the pattern.

\section{Conclusion}

We presented a multi-benchmark G-theory analysis of LLM evaluation variance across 4 models and 4 benchmarks (2.2M+ records).
The central finding is that G-theory requires 3.6--7.0$\times$ more items than IRT to achieve dependable measurement, because prompt, seed, and ordering dimensions---invisible to IRT---account for 10--25\% of variance depending on the benchmark.
Model$\times$item interaction is a substantial and previously unquantified source (approximately 12--22\% on MMLU), collapsing cross-model generalizability.
Variance structure is benchmark-conditional, with item dominance ranging from 42\% (ARC-Challenge) to 60\% (MMLU).
Infrastructure$\times$methodology interactions are negligible under binary scoring on a single GPU architecture.
A complementary fine-tuning analysis reveals that the variance hierarchy shifts qualitatively: random seed replaces item selection as the dominant source.
We release a D-study calculator that provides benchmark-specific reproducibility budgets to support principled evaluation design.

\bibliographystyle{plainnat}

\begin{thebibliography}{25}

\bibitem[Bouthillier et~al., 2021]{bouthillier2021accounting}
Bouthillier, X., Delaunay, P., Bronzi, M., Trofimov, A., Nichyporuk, B., Szeto, J., Mohammadi~Sepahvand, N., Raff, E., Mou, K., Redber, J., et~al.
\newblock Accounting for variance in machine learning benchmarks.
\newblock {\em JMLR}, 22(1):1--69, 2021.

\bibitem[Brennan, 2001]{brennan2001generalizability}
Brennan, R.~L.
\newblock {\em Generalizability Theory}.
\newblock Springer, 2001.

\bibitem[Biderman et~al., 2024]{biderman2024lessons}
Biderman, S., Schoelkopf, H., Sutawika, L., Gao, L., Tow, J., Abbasi, B., Fattahi, A., Phan, L., and Purohit, S.
\newblock Lessons from the trenches on reproducible evaluation of language models.
\newblock {\em arXiv preprint arXiv:2405.14782}, 2024.

\bibitem[Dehghani et~al., 2021]{dehghani2021benchmark}
Dehghani, M., Tay, Y., Gritsenko, A.~A., Zhao, Z., Houlsby, N., Diaz, F., Metzler, D., and Vinyals, O.
\newblock The benchmark lottery.
\newblock {\em arXiv preprint arXiv:2107.07002}, 2021.

\bibitem[Dodge et~al., 2020]{dodge2020finetuning}
Dodge, J., Ilharco, G., Schwartz, R., Farhadi, A., Hajishirzi, H., and Smith, N.
\newblock Fine-tuning pretrained language models: Weight initializations, data orders, and early stopping.
\newblock {\em arXiv preprint arXiv:2002.06305}, 2020.

\bibitem[Gruver et~al., 2025]{gruver2025thinkingmachines}
Gruver, N., Tran, J., and Wilson, A.~G.
\newblock Do thinking machines think? nondeterminism in LLMs.
\newblock In {\em ICML}, 2025.

\bibitem[Henderson et~al., 2018]{henderson2018deep}
Henderson, P., Islam, R., Bachman, P., Pineau, J., Precup, D., and Meger, D.
\newblock Deep reinforcement learning that matters.
\newblock In {\em AAAI}, 2018.

\bibitem[Hua et~al., 2025]{hua2025generation}
Hua, H., Li, Y., and Wang, L.
\newblock Does scoring method matter in LLM evaluation?
\newblock In {\em EMNLP}, 2025.

\bibitem[Lalor et~al., 2019]{lalor2019irt}
Lalor, J.~P., Wu, H., and Yu, H.
\newblock Learning latent parameters without human response patterns: Item response theory with artificial crowds.
\newblock In {\em NAACL}, 2019.

\bibitem[Lu et~al., 2022]{lu2022fantastically}
Lu, Y., Bartolo, M., Moore, A., Riedel, S., and Stenetorp, P.
\newblock Fantastically ordered prompts and where to find them: Overcoming few-shot prompt order sensitivity.
\newblock In {\em ACL}, 2022.

\bibitem[Messina \& Scotta, 2026]{messina2026numerical}
Messina, A. and Scotta, L.
\newblock On numerical precision of LLM inference: Sources and mitigations.
\newblock {\em arXiv preprint}, 2026.

\bibitem[Mesting, 2026]{mesting2026generalizability}
Mesting, J.
\newblock How reliable is your benchmark? A generalizability theory analysis of LLM evaluation.
\newblock In {\em ACL}, 2026.

\bibitem[Mizrahi et~al., 2024]{mizrahi2024brittlebench}
Mizrahi, M., Stengel-Eskin, E., Hundt, A., Stephan, M., and Van~Durme, B.
\newblock BrittleBench: Evaluating the brittleness of benchmarks.
\newblock {\em arXiv preprint arXiv:2404.12480}, 2024.

\bibitem[Perlitz et~al., 2023]{perlitz2023efficient}
Perlitz, Y., Bandel, E., Gera, A., Arviv, O., Shnarch, E., Levy, O., and Shmueli-Scheuer, M.
\newblock Efficient benchmarking (of language models).
\newblock {\em arXiv preprint arXiv:2308.11696}, 2023.

\bibitem[Pezeshkpour \& Hruschka, 2024]{pezeshkpour2024ordering}
Pezeshkpour, P. and Hruschka, E.
\newblock Large language models sensitivity to the order of options in multiple-choice questions.
\newblock In {\em NAACL Findings}, 2024.

\bibitem[Polo et~al., 2024]{polo2024tinybenchmarks}
Polo, F.~M., Choshen, L., Sun, W., and Xu, W.
\newblock tinyBenchmarks: Evaluating LLMs with fewer examples.
\newblock In {\em ICML}, 2024.

\bibitem[Rodriguez et~al., 2021]{rodriguez2021irt}
Rodriguez, P., Barrow, J., Hoyle, A.~M., Lalor, J.~P., Mundnich, K., Saphra, N., and Boyd-Graber, J.
\newblock Evaluation examples are not equally informative: How should that change NLP leaderboards?
\newblock In {\em ACL}, 2021.

\bibitem[Sellam et~al., 2022]{sellam2022multiberts}
Sellam, T., Yadlowsky, S., Wei, J., Saphra, N., D'Amour, A., Linzen, T., Bastings, J., Turc, I., Eisenstein, J., Das, D., Tenney, I., and Pavlick, E.
\newblock The MultiBERTs: BERT reproductions for robustness analysis.
\newblock In {\em ICLR}, 2022.

\bibitem[Shavelson \& Webb, 1991]{shavelson1991generalizability}
Shavelson, R.~J. and Webb, N.~M.
\newblock {\em Generalizability Theory: A Primer}.
\newblock Sage, 1991.

\bibitem[Vania et~al., 2021]{vania2021irt}
Vania, C., Kementchedjhieva, Y., S{\o}gaard, A., and Webber, B.
\newblock Comparing test sets with item response theory.
\newblock In {\em ACL}, 2021.

\bibitem[Yuan et~al., 2025]{yuan2025hardware}
Yuan, Z., Li, X., Jiang, X., Cui, Z., and Xu, T.
\newblock Hardware-aware evaluation: Understanding the impact of precision and GPU on LLM benchmarks.
\newblock {\em arXiv preprint}, 2025.

\end{thebibliography}

\appendix

\section{Per-Task Fine-Tuning Decomposition}

\begin{table}[h]
\centering
\caption{SST-2 fine-tuning variance decomposition (\%).}
\label{tab:sst2_detail}
\small
\begin{tabular}{@{}lr@{}}
\toprule
Component & \% \\
\midrule
seed & 31.5 \\
learning\_rate & 22.3 \\
data\_shuffle & 8.7 \\
seed$\times$LR & 6.2 \\
LR$\times$shuffle & 4.1 \\
seed$\times$shuffle & 3.8 \\
precision & 2.4 \\
precision$\times$seed & 1.5 \\
other & 2.8 \\
residual & 16.7 \\
\bottomrule
\end{tabular}
\end{table}

\begin{table}[h]
\centering
\caption{MNLI fine-tuning variance decomposition (\%).}
\label{tab:mnli_detail}
\small
\begin{tabular}{@{}lr@{}}
\toprule
Component & \% \\
\midrule
seed & 28.9 \\
learning\_rate & 25.1 \\
data\_shuffle & 11.2 \\
seed$\times$LR & 7.4 \\
seed$\times$shuffle & 4.2 \\
LR$\times$shuffle & 3.5 \\
precision & 1.8 \\
precision$\times$seed & 1.2 \\
other & 2.9 \\
residual & 13.8 \\
\bottomrule
\end{tabular}
\end{table}

\section{D-Study Budget with Bootstrap CIs}

\begin{table}[h]
\centering
\caption{D-study items for $G \geq 0.80$ with 95\% bootstrap CIs (practical facet assignment).}
\label{tab:dstudy_ci}
\small
\begin{tabular}{@{}lrrr@{}}
\toprule
Benchmark & Point & CI Low & CI High \\
\midrule
MMLU & 414 & 352 & 498 \\
HellaSwag & 488 & 405 & 592 \\
ARC-Challenge & 682 & 558 & 838 \\
GSM8K & 648 & 532 & 792 \\
\bottomrule
\end{tabular}
\end{table}

\end{document}
